\chapter{2010年山东卷（理）}

\section{单选题}

\begin{problem}
已知全集 \(U = \R\)，集合 \(M = \{x \mid |x - 1| \leqslant 2\}\)，则 \(\complement_U M =\)（\quad）

\choices
    {\( \{x \mid -1 < x < 3\} \)}
    {\(\{x\mid -1\leqslant x\leqslant 3\}\)}
    {\(\{x \mid x < -1\text{ 或 }x > 3\}\)}
    {\(\{x\mid x\leqslant -1\text{ 或 }x\geqslant 3\}\)}
\end{problem}

\begin{answer}
C
\end{answer}

\begin{solution}
由 \(|x-1|\leqslant2\)，得 \(-2\leqslant x-1\leqslant2\)，即 \(-1\leqslant x\leqslant3\). 因此 \(M=[-1,3]\). 由于全集为 \(\R\)，所以 \(\complement_U M=(-\infty,-1)\cup(3,+\infty)\)，故选 C.
\end{solution}

\begin{problem}
已知 \(\frac{a + 2\i}{\i} = b + \i\) （\(a\)，\(b \in \R\)），其中 \(\i\) 为虚数单位，则 \(a + b =\)（\quad）

\choices
    {\(-1\)}
    {\(1\)}
    {\(2\)}
    {\(3\)}
\end{problem}

\begin{answer}
B
\end{answer}

\begin{solution}
因为 \(\dfrac{1}{\i}=-\i\)，所以
\[
\frac{a+2\i}{\i}=-a\i+2
\]
与 \(b+\i\) 比较实部和虚部，得 \(b=2\)，\(-a=1\)，即 \(a=-1\). 因而 \(a+b=1\)，故选 B.
\end{solution}

\begin{problem}
在空间，下列命题正确的是（\quad）

\choices
    {平行直线的平行投影重合}
    {平行于同一直线的两个平面平行}
    {垂直于同一平面的两个平面平行}
    {垂直于同一平面的两条直线平行}
\end{problem}

\begin{answer}
D
\end{answer}

\begin{solution}
选项 A 不正确，因为两条平行直线的平行投影可以是两条不同的平行直线，不一定重合；选项 B 不正确，因为平行于同一直线的两个平面可能相交；选项 C 不正确，因为垂直于同一平面的两个平面可能相交.

若两条直线都垂直于同一平面，则它们都与该平面的法线方向平行，所以这两条直线互相平行，选项 D 正确. 故选 D.
\end{solution}

\begin{problem}
设 \( f(x) \) 为定义在 \(\R\) 上的奇函数，当 \( x \geqslant 0 \) 时， \( f(x) = 2^{x} + 2x + b \) （\(b\) 为常数），则 \( f(-1) = \)（\quad）

\choices
    {\(3\)}
    {\(1\)}
    {\(-1\)}
    {\(-3\)}
\end{problem}

\begin{answer}
D
\end{answer}

\begin{solution}
将 \(x=0\) 代入题设表达式，得 \(f(0)=1+b\). 奇函数满足 \(f(0)=0\)，所以 \(b=-1\). 因而 \(f(1)=2^1+2-1=3\). 由奇函数性质，\(f(-1)=-f(1)=-3\)，故选 D.
\end{solution}

\begin{problem}
已知随机变量 \(\xi\) 服从正态分布 \(N(0, \sigma^2)\). 若 \(P(\xi > 2) = 0.023\)，则 \(P(-2 \leqslant \xi \leqslant 2) =\)（\quad）

\choices
    {\(0.477\)}
    {\(0.628\)}
    {\(0.954\)}
    {\(0.977\)}
\end{problem}

\begin{answer}
C
\end{answer}

\begin{solution}
正态分布 \(N(0,\sigma^2)\) 的图象关于纵轴对称，因此 \(P(\xi<-2)=P(\xi>2)=0.023\). 于是
\[
P(-2\leqslant\xi\leqslant2)=1-P(\xi<-2)-P(\xi>2)=1-2\times0.023=0.954
\]
故选 C.
\end{solution}

\begin{problem}
样本中共有五个个体，其值分别为 \(a\)，\(0\)，\(1\)，\(2\)，\(3\). 若该样本的平均值为 1，则样本方差为（\quad）

\choices
    {\(\sqrt{\frac{6}{5}}\)}
    {\(\frac{6}{5}\)}
    {\(\sqrt{2}\)}
    {\(2\)}
\end{problem}

\begin{answer}
D
\end{answer}

\begin{solution}
由平均值为 \(1\)，得
\[
\frac{a+0+1+2+3}{5}=1
\]
所以 \(a=-1\). 此时样本方差为
\[
s^2=\frac{(-1-1)^2+(0-1)^2+(1-1)^2+(2-1)^2+(3-1)^2}{5}=\frac{4+1+0+1+4}{5}=2
\]
故选 D.
\end{solution}

\begin{problem}
由曲线 \( y = x^2 \)，\( y = x^3 \) 围成的封闭图形面积为（\quad）

\choices
    {\(\frac{1}{12}\)}
    {\(\frac{1}{4}\)}
    {\(\frac{1}{3}\)}
    {\(\frac{7}{12}\)}
\end{problem}

\begin{answer}
A
\end{answer}

\begin{solution}
两条曲线的交点满足 \(x^2=x^3\)，即 \(x=0\) 或 \(x=1\). 在 \([0,1]\) 上有 \(x^2\geqslant x^3\)，所以封闭图形面积为
\[
\int_0^1(x^2-x^3)\,\mathrm{d}x=\left.\left(\frac{x^3}{3}-\frac{x^4}{4}\right)\right|_0^1=\frac{1}{12}
\]
故选 A.
\end{solution}

\begin{problem}
某台小型晚会由 6 个节目组成. 演出顺序有如下要求：节目甲必须排在前两位，节目乙不能排在第一位，节目丙必须排在最后一位，该台晚会节目演出顺序的编排方案共有（\quad）

\choices
    {\(36\text{ 种}\)}
    {\(42\text{ 种}\)}
    {\(48\text{ 种}\)}
    {\(54\text{ 种}\)}
\end{problem}

\begin{answer}
B
\end{answer}

\begin{solution}
节目丙固定在第 6 位，只需安排节目甲、节目乙及另外 3 个节目在前 5 位.

若节目甲在第 1 位，节目乙可在第 2 至第 5 位，有 \(4\) 种选法，其余 3 个节目有 \(3!\) 种排法，共 \(4\times3!=24\) 种.

若节目甲在第 2 位，第 1 位不能安排节目乙，只能从另外 3 个节目中选，有 \(3\) 种选法；剩余 3 个位置安排节目乙及另外 2 个节目，有 \(3!\) 种排法，共 \(3\times3!=18\) 种.

所以编排方案共有 \(24+18=42\) 种，故选 B.
\end{solution}

\begin{problem}
设 \(\{a_{n}\}\) 是等比数列，则 \(a_{1} < a_{2} < a_{3}\) 是“数列 \(\{a_{n}\}\) 是递增数列”的（\quad）

\choices
    {充分而不必要条件}
    {必要而不充分条件}
    {充分必要条件}
    {既不充分也不必要条件}
\end{problem}

\begin{answer}
C
\end{answer}

\begin{solution}
设等比数列的首项为 \(a_1\)，公比为 \(q\). 若 \(a_1<a_2<a_3\)，则
\(a_1(q-1)>0\)，\(a_1q(q-1)>0\).
两式相除得 \(q>0\). 再结合第一式可知：当 \(q>1\) 时 \(a_1>0\)；当 \(0<q<1\) 时 \(a_1<0\). 在这两种情形下，对任意 \(n\in\mathbf N^*\)，都有
\[
a_{n+1}-a_n=a_1q^{n-1}(q-1)>0
\]
所以数列递增.

反过来，若数列递增，则必有 \(a_1<a_2<a_3\). 因而这两个条件互为充分必要条件，故选 C.
\end{solution}

\begin{problem}
设变量 \(x\)，\(y\) 满足约束条件 \(\begin{cases}
x - y + 2 \geqslant 0,\\
x - 5y + 10 \leqslant 0,\\
x + y - 8 \leqslant 0,
\end{cases}\) 则目标函数 \(z = 3x - 4y\) 的最大值和最小值分别为（\quad）

\choices
    {\(3\)，\(-11\)}
    {\(-3\)，\(-11\)}
    {\(11\)，\(-3\)}
    {\(11\)，\(3\)}
\end{problem}

\begin{answer}
A
\end{answer}

\begin{solution}
由约束条件得
\(y\leqslant x+2\)，\(y\geqslant\frac{x}{5}+2\)，\(y\leqslant8-x\).
三条边界直线的交点分别由
\(y=x+2\)，\(y=\frac{x}{5}+2;\)，\(y=x+2\)，\(y=8-x;\)，\(y=\frac{x}{5}+2\)，\(y=8-x\)
求得为 \((0,2)\)、\((3,5)\)、\((5,3)\). 这三个交点都满足约束条件，且可行域为它们围成的三角形. 在三个顶点处计算目标函数 \(z=3x-4y\)，得到
\(z(0,2)=-8\)，\(z(3,5)=-11\)，\(z(5,3)=3\).
所以最大值为 \(3\)，最小值为 \(-11\)，故选 A.
\end{solution}

\begin{problem}
函数 \( y = 2^{x} - x^{2} \) 的图象大致是（\quad）

\choices
    {\choicebitmap[width=0.15\paperwidth]{img/2010/shandong_science/q11_choiceA.png}}
    {\choicebitmap[width=0.15\paperwidth]{img/2010/shandong_science/q11_choiceB.png}}
    {\choicebitmap[width=0.15\paperwidth]{img/2010/shandong_science/q11_choiceC.png}}
    {\choicebitmap[width=0.15\paperwidth]{img/2010/shandong_science/q11_choiceD.png}}
\end{problem}

\begin{answer}
A
\end{answer}

\begin{solution}
令 \(f(x)=2^x-x^2\)，则
\(f'(x)=2^x\ln2-2x\)，\(f''(x)=2^x(\ln2)^2-2\).
当 \(x\leqslant0\) 时，\(f'(x)>0\)，所以函数在负半轴上递增. 对正半轴，\(f''(x)\) 严格递增，故 \(f'(x)\) 先减后增；又
\(f'(0)=\ln2>0\)，\(f'(1)=2\ln2-2<0\)，\(f'(4)=16\ln2-8>0\).
因此函数图象先递增，在正半轴上先出现一个极大值，再递减出现一个极小值，最后递增. 另外 \(f(0)=1>0\)，且当 \(x\to-\infty\) 时 \(f(x)\to-\infty\)，当 \(x\to+\infty\) 时 \(f(x)\to+\infty\). 进一步，\(f(1)=1>0\)、\(f(2)=0\)、\(f(3)=-1<0\)、\(f(4)=0\)，说明图象在正半轴上经过 \((2,0)\)、\((4,0)\)，并在第二个极值点附近先位于 \(x\) 轴下方再上升. 因而只能选 A.
\end{solution}

\begin{problem}
定义平面向量之间的一种运算“\(\odot\)”如下：对任意的 \(\bs{a} = (m, n)\)，\(\bs{b} = (p, q)\). 令 \(\bs{a} \odot \bs{b} = mq - np\). 下面说法错误的是（\quad）

\choices
    {若 \(\bs{a}\) 与 \(\bs{b}\) 共线，则 \(\bs{a} \odot \bs{b} = 0\)}
    {\(\bs{a} \odot \bs{b} = \bs{b} \odot \bs{a}\)}
    {对任意的 \(\lambda \in \R\)，有 \((\lambda \bs{a}) \odot \bs{b} = \lambda (\bs{a} \odot \bs{b})\)}
    {\((\bs{a}\odot \bs{b})^2 + (\bs{a}\cdot \bs{b})^2 = |\bs{a}|^2 |\bs{b}|^2\)}
\end{problem}

\begin{answer}
B
\end{answer}

\begin{solution}
由定义，\(\bs{a}\odot\bs{b}=mq-np\).

若 \(\bs{a}\) 与 \(\bs{b}\) 共线，则 \(mq-np=0\)，所以选项 A 正确. 交换两向量得
\[
\bs{b}\odot\bs{a}=pn-qm=-(mq-np)=-\bs{a}\odot\bs{b}
\]
一般不等于 \(\bs{a}\odot\bs{b}\)，例如取 \(\bs{a}=(1,0)\)，\(\bs{b}=(0,1)\)，则 \(\bs{a}\odot\bs{b}=1\)，而 \(\bs{b}\odot\bs{a}=-1\)，所以选项 B 错误.

又
\[
(\lambda\bs{a})\odot\bs{b}=(\lambda m)q-(\lambda n)p=\lambda(mq-np)=\lambda(\bs{a}\odot\bs{b})
\]
所以选项 C 正确. 对选项 D，直接展开左端得
\[
\begin{aligned}
(mq-np)^2+(mp+nq)^2
&=(m^2q^2-2mnpq+n^2p^2) \\
&\quad +(m^2p^2+2mnpq+n^2q^2) \\
&=(m^2+n^2)(p^2+q^2)=|\bs{a}|^2|\bs{b}|^2,
\end{aligned}
\]
即选项 D 正确. 故错误的是 B.
\end{solution}

\section{填空题}

\begin{problem}
执行如图所示的程序框图，若输入 \(x = 10\)，则输出 \(y\) 的值为\fillinblank{}.

\texfigure[max width=0.36\linewidth]{img_repaint/2010/shandong_science/q13_fig1.tex}
\end{problem}

\begin{answer}
\(-\frac{5}{4}\)
\end{answer}

\begin{solution}
依次迭代计算：初始输入 \(x=10\) 时，\(y=\frac12x-1=4\)，且 \(|y-x|=6\geqslant1\)，令 \(x=y=4\)；再得 \(y=1\)，且 \(|1-4|=3\geqslant1\)，令 \(x=1\)；再得 \(y=-\frac12\)，且 \(\left|-\frac12-1\right|=\frac32\geqslant1\)，令 \(x=-\frac12\)；最后得 \(y=-\frac54\)，且
\[
\left|y-x\right|=\left|-\frac54+\frac12\right|=\frac34<1
\]
此时满足判断条件，输出 \(y=-\frac54\).
\end{solution}

\begin{problem}
若对任意 \(x > 0\)，\(\frac{x}{x^2 + 3x + 1} \leqslant a\) 恒成立，则 \(a\) 的取值范围是\fillinblank{}.
\end{problem}

\begin{answer}
\(a \geqslant \frac{1}{5}\)
\end{answer}

\begin{solution}
因 \(x>0\)，故 \(x^2+3x+1>0\). 注意到
\[
x^2+3x+1-5x=(x-1)^2\geqslant0
\]
所以 \(x^2+3x+1\geqslant5x\)，从而
\[
\frac{x}{x^2+3x+1}\leqslant\frac15
\]
当 \(x=1\) 时等号成立，因此左端的最大值为 \(\frac15\). 题设不等式对任意 \(x>0\) 恒成立，当且仅当 \(a\) 不小于该最大值，即 \(a\geqslant\frac15\).
\end{solution}

\begin{problem}
在 \(\triangle ABC\) 中，\(A\)，\(B\)，\(C\) 所对的边分别为 \(a\)，\(b\)，\(c\). 若 \(a = \sqrt{2}\)，\(b = 2\)，\(\sin B + \cos B = \sqrt{2}\)，则角 \(A\) 的大小为\fillinblank{}.
\end{problem}

\begin{answer}
\(\frac{\pi}{6}\)
\end{answer}

\begin{solution}
因为 \(0<B<\pi\)，所以
\[
\sin B+\cos B=\sqrt2\sin\left(B+\frac\pi4\right)=\sqrt2
\]
又 \(B+\frac\pi4\in\left(\frac\pi4,\frac{5\pi}{4}\right)\)，故只能有 \(B+\frac\pi4=\frac\pi2\)，从而 \(B=\frac\pi4\).

由正弦定理，
\[
\frac{a}{\sin A}=\frac{b}{\sin B}
\]
所以
\[
\sin A=\frac{a\sin B}{b}=\frac{\sqrt2\cdot\frac{\sqrt2}{2}}{2}=\frac12
\]
三角形内角满足 \(A<\pi-B=\frac{3\pi}{4}\)，因此 \(A\) 不可能为 \(\frac{5\pi}{6}\)，只能为 \(\frac\pi6\).
\end{solution}

\begin{problem}
已知圆 \(C\) 过点 \((1,0)\)，且圆心在 \(x\) 轴的正半轴上，直线 \(l: y = x - 1\) 被圆 \(C\) 所截得的弦长为 \(2\sqrt{2}\)，则过圆心且与直线 \(l\) 垂直的直线的方程为\fillinblank{}.
\end{problem}

\begin{answer}
\(x+y-3=0\)
\end{answer}

\begin{solution}
设圆心为 \(O(r,0)\)，其中 \(r>0\). 因圆过点 \((1,0)\)，圆的半径为 \(R=|r-1|\). 直线 \(l:y=x-1\) 可写为 \(x-y-1=0\)，所以圆心到直线 \(l\) 的距离为
\[
d=\frac{|r-1|}{\sqrt2}
\]
由弦长公式，
\[
2\sqrt{R^2-d^2}=2\sqrt2
\]
于是
\[
(r-1)^2-\frac{(r-1)^2}{2}=2
\]
即 \((r-1)^2=4\). 结合 \(r>0\)，得 \(r=3\)，故圆心为 \((3,0)\).

直线 \(l\) 的斜率为 \(1\)，过圆心且垂直于 \(l\) 的直线斜率为 \(-1\)，其方程为 \(y=-(x-3)\)，即 \(x+y-3=0\).
\end{solution}

\section{解答题}

\begin{problem}
已知函数 \( f(x) = \frac{1}{2} \sin 2x \sin \varphi + \cos^2 x \cos \varphi - \frac{1}{2} \sin \left( \frac{\pi}{2} + \varphi \right) \) （\(0 < \varphi < \pi\)），其图象过点 \( \left( \frac{\pi}{6}, \frac{1}{2} \right) \).
\begin{enumerate}
\item 求 \(\varphi\) 的值；
\item 将函数 \( y = f(x) \) 的图象上各点的横坐标缩短到原来的 \( \frac{1}{2} \)，纵坐标不变，得到函数 \( y = g(x) \) 的图象. 求函数 \( g(x) \) 在 \( \left[0, \frac{\pi}{4}\right] \) 上的最大值和最小值.
\end{enumerate}
\end{problem}


\begin{solution}
\begin{enumerate}
\item 由 \(\cos^2x=\frac{1+\cos2x}{2}\)，且 \(\sin\left(\frac{\pi}{2}+\varphi\right)=\cos\varphi\)，得
\[
f(x)=\frac12\sin2x\sin\varphi+\frac12\cos2x\cos\varphi=\frac12\cos\left(2x-\varphi\right)
\]
因函数图象过点 \(\left(\frac{\pi}{6},\frac12\right)\)，所以
\[
\frac12\cos\left(\frac{\pi}{3}-\varphi\right)=\frac12
\]
即 \(\cos\left(\frac{\pi}{3}-\varphi\right)=1\). 由 \(0<\varphi<\pi\)，可知 \(-\frac{2\pi}{3}<\frac{\pi}{3}-\varphi<\frac{\pi}{3}\)，故只能有 \(\frac{\pi}{3}-\varphi=0\)，从而 \(\varphi=\frac{\pi}{3}\).
\item 横坐标缩短到原来的 \(\frac12\) 后，横坐标为 \(x\) 的点对应原图象上横坐标为 \(2x\) 的点，因此
\[
g(x)=f(2x)=\frac12\cos\left(4x-\frac{\pi}{3}\right)
\]
当 \(x\in\left[0,\frac{\pi}{4}\right]\) 时，令 \(t=4x-\frac{\pi}{3}\)，则 \(t\in\left[-\frac{\pi}{3},\frac{2\pi}{3}\right]\). 在此区间内，\(\cos t\) 的最大值为 \(1\)，最小值为 \(\cos\frac{2\pi}{3}=-\frac12\)，所以
\(g(x)_{\max}=\frac12\)，\(g(x)_{\min}=-\frac14\).
\end{enumerate}
\end{solution}

\begin{problem}
已知等差数列 \(\{a_{n}\}\) 满足：\(a_{3} = 7\)，\(a_{5} + a_{7} = 26\). 数列 \(\{a_{n}\}\) 的前 \(n\) 项和为 \(S_{n}\).
\begin{enumerate}
\item 求 \(a_{n}\) 及 \(S_{n}\)；
\item 令 \( b_{n} = \frac{1}{a_{n}^{2} - 1} \)（\(n \in \mathbf{N}^{*}\)），求数列 \(\{b_{n}\}\) 的前 \(n\) 项和 \(T_{n}\).
\end{enumerate}
\end{problem}

\begin{solution}
\begin{enumerate}
\item 设等差数列 \(\{a_n\}\) 的公差为 \(d\)，则
\(a_1+2d=7\)，\((a_1+4d)+(a_1+6d)=26\).
由第一式得 \(a_1=7-2d\)，代入第二式得
\[
2(7-2d)+10d=26
\]
所以 \(d=2\)，\(a_1=3\). 因而
\(a_n=a_1+(n-1)d=2n+1\)，\(S_n=\frac{n(a_1+a_n)}2=n(n+2)\).
\item 由 \(a_n=2n+1\)，得
\[
b_n=\frac{1}{(2n+1)^2-1}=\frac{1}{4n(n+1)}=\frac14\left(\frac1n-\frac{1}{n+1}\right)
\]
所以
\[
T_n=\sum_{k=1}^{n}b_k=\frac14\sum_{k=1}^{n}\left(\frac1k-\frac{1}{k+1}\right)=\frac14\left(1-\frac{1}{n+1}\right)=\frac{n}{4(n+1)}
\]
\end{enumerate}
\end{solution}

\begin{problem}
如图，在五棱锥 \(P - ABCDE\) 中，\(PA \perp\) 平面 \(ABCDE\)，\(AB \parallel CD\)，\(AC \parallel ED\)，\(AE \parallel BC\)，\(\angle ABC = 45^\circ\)，\(AB = 2\sqrt{2}\)，\(BC = 2AE = 4\)，三角形 \(PAB\) 是等腰三角形.
\begin{enumerate}
\item 求证：平面 \(PCD \perp\) 平面 \(PAC\)；
\item 求直线 \(PB\) 与平面 \(PCD\) 所成角的大小；
\item 求四棱锥 \(P - ACDE\) 的体积.
\end{enumerate}

\bitmapfigure[max width=0.34\linewidth]{img/2010/shandong_science/q19_fig1.png}
\end{problem}

\begin{solution}
\begin{enumerate}
\item 在三角形 \(ABC\) 中，由余弦定理得
\[
AC^2=AB^2+BC^2-2AB\cdot BC\cos\angle ABC=8+16-2\cdot2\sqrt2\cdot4\cdot\frac{\sqrt2}{2}=8
\]
于是 \(AB^2+AC^2=8+8=16=BC^2\)，所以 \(\angle BAC=90^\circ\)，即 \(AB\perp AC\). 又因为 \(CD\parallel AB\)，故 \(CD\perp AC\). 由于 \(PA\perp\) 平面 \(ABCDE\)，而 \(CD\) 在该平面内，所以 \(CD\perp PA\). 因此 \(CD\) 垂直于平面 \(PAC\)，而 \(CD\subset\) 平面 \(PCD\)，从而平面 \(PCD\perp\) 平面 \(PAC\).
\item 取直角坐标系，使平面 \(ABCDE\) 为 \(z=0\)，取
\(A(0,0,0)\)，\(B(-2,-2,0)\)，\(C(2,-2,0)\).
因为三角形 \(PAB\) 是直角三角形，且直角在 \(A\)，所以等腰的两腰为 \(PA\)，\(AB\)，从而 \(PA=AB=2\sqrt2\)，可取 \(P(0,0,2\sqrt2)\). 又因为 \(AE\parallel BC\)，且 \(AE=2\)，按五边形的顺序取 \(E(2,0,0)\).

设 \(D=C+\mu(B-A)=\left(2-2\mu,-2-2\mu,0\right)\). 由 \(ED\parallel AC\)，设 \(D=E+\nu(C-A)=\left(2+2\nu,-2\nu,0\right)\). 比较坐标得 \(\nu=-\mu\)，且 \(\nu=1+\mu\)，所以 \(\mu=-\frac12\)，\(\nu=\frac12\)，即 \(D(3,-1,0)\).

于是
\(\overrightarrow{CD}=(1,1,0)\)，\(\overrightarrow{PC}=(2,-2,-2\sqrt2)\)，\(\overrightarrow{PB}=(-2,-2,-2\sqrt2)\).
取平面 \(PCD\) 的一个法向量
\[
\bs{n}=\overrightarrow{CD}\times\overrightarrow{PC}=(-2\sqrt2,2\sqrt2,-4)
\]
设直线 \(PB\) 与平面 \(PCD\) 所成的角为 \(\theta\)，则
\[
\sin\theta=\frac{\left|\overrightarrow{PB}\cdot\bs{n}\right|}{\left|\overrightarrow{PB}\right|\left|\bs{n}\right|}=\frac{8\sqrt2}{4\cdot4\sqrt2}=\frac12
\]
因为 \(0\leqslant\theta\leqslant\frac\pi2\)，所以 \(\theta=\frac\pi6\).
\item 由上述坐标，四边形 \(ACDE\) 可分成三角形 \(ACD\) 和三角形 \(ADE\)，故
\(S_{\triangle ACD}=\frac12\left|2\cdot(-1)-(-2)\cdot3\right|=2\)，\(S_{\triangle ADE}=\frac12\left|3\cdot0-(-1)\cdot2\right|=1\).
所以 \(S_{ACDE}=3\). 四棱锥 \(P-ACDE\) 的高为 \(PA=2\sqrt2\)，从而
\[
V_{P-ACDE}=\frac13S_{ACDE}\cdot PA=\frac13\cdot3\cdot2\sqrt2=2\sqrt2
\]
\end{enumerate}
\end{solution}

\begin{problem}
某学校举行知识竞赛，第一轮选拔共设有 A，B，C，D 四个问题，规则如下：
\begin{enumerate}
\item 每位参加者计分器的初始分均为 10 分，答对问题 A，B，C，D 分别加 1 分，2 分，3 分，6 分，答错任一题减 2 分；
\item 每回答一题，计分器显示累计分数，当累计分数小于 8 分时，答题结束，淘汰出局；当累计分数大于或等于 14 分时，答题结束，进入下一轮；当答完四题，累计分数仍不足 14 分时，答题结束，淘汰出局；
\item 每位参加者按问题 A，B，C，D 顺序作答，直至答题结束.
\end{enumerate}
假设甲同学对问题 A，B，C，D 回答正确的概率依次为 \(\frac{3}{4}\)，\(\frac{1}{2}\)，\(\frac{1}{3}\)，\(\frac{1}{4}\)，且各题回答正确与否相互之间没有影响.
\begin{enumerate}
\item 求甲同学能进入下一轮的概率；
\item 用 \(\xi\) 表示甲同学本轮答题结束时答题的个数，求 \(\xi\) 的分布列和数学期望 \(E(\xi)\).
\end{enumerate}
\end{problem}

\begin{solution}
\begin{enumerate}
\item 设第 \(i\) 个问题答对为事件 \(M_i\)，答错为事件 \(N_i\)，则
\(P(M_1)=\frac34\)，\(P(M_2)=\frac12\)，\(P(M_3)=\frac13\)，\(P(M_4)=\frac14\)
且各事件相互独立. 答完第一个问题后，累计分数为 \(11\) 或 \(8\)，均可继续答题. 答完第二个问题后，只有发生 \(N_1N_2\) 时累计分数为 \(6<8\)，会停止答题；其他情形均可继续.

若前三题答题结果为 \(M_1M_2M_3\)，前三题得分为 \(16\)，此时进入下一轮. 到达第四题的四种前三题结果及分数分别为 \(M_1M_2N_3\)、\(M_1N_2M_3\)、\(N_1M_2M_3\)、\(N_1M_2N_3\)，对应分数为 \(11\)、\(12\)、\(13\)、\(8\)，这四种情形只有答对第四题才能进入下一轮. 结果为 \(M_1N_2N_3\) 时，前三题得分为 \(7<8\)，已经在第三题后淘汰，不能继续回答第四题. 因而进入下一轮的概率为
\[
\begin{aligned}
P&=P(M_1M_2M_3)+P(M_1M_2N_3M_4)+P(M_1N_2M_3M_4)\\
&\quad+P(N_1M_2M_3M_4)+P(N_1M_2N_3M_4)\\
&=\frac18+\frac1{16}+\frac1{32}+\frac1{96}+\frac1{48}=\frac14.
\end{aligned}
\]
\item 随机变量 \(\xi\) 的可能取值为 \(2\)，\(3\)，\(4\). 其中
\[
P(\xi=2)=P(N_1N_2)=\frac14\cdot\frac12=\frac18
\]
\[
P(\xi=3)=P(M_1M_2M_3)+P(M_1N_2N_3)
=\frac18+\frac14=\frac38
\]
\[
P(\xi=4)=P(M_1M_2N_3)+P(M_1N_2M_3)+P(N_1M_2M_3)+P(N_1M_2N_3)
=\frac14+\frac18+\frac1{24}+\frac1{12}=\frac12
\]
所以 \(\xi\) 的分布列为
\begin{center}
\begin{tabular}{c|ccc}
\(\xi\) & \(2\) & \(3\) & \(4\) \\
\hline
\(P\) & \(\frac18\) & \(\frac38\) & \(\frac12\)
\end{tabular}
\end{center}
数学期望为
\[
E(\xi)=2\cdot\frac18+3\cdot\frac38+4\cdot\frac12=\frac{27}{8}
\]
\end{enumerate}
\end{solution}

\begin{problem}
如图，已知椭圆 \(\frac{x^2}{a^2} + \frac{y^2}{b^2} = 1\) （\(a > b > 0\)）的离心率为 \(\frac{\sqrt{2}}{2}\)，以该椭圆上的点和椭圆的左，右焦点 \(F_1\)，\(F_2\) 为顶点的三角形的周长为 \(4(\sqrt{2} + 1)\). 一等轴双曲线的顶点是该椭圆的焦点. 设 \(P\) 为该双曲线上异于顶点的任一点，直线 \(PF_1\) 和 \(PF_2\) 与椭圆的交点分别为 \(A\)，\(B\) 和 \(C\)，\(D\).
\begin{enumerate}
\item 求椭圆和双曲线的标准方程；
\item 设直线 \(PF_{1}\)，\(PF_{2}\) 的斜率分别为 \(k_{1}\)，\(k_{2}\)，证明：\(k_{1} \cdot k_{2} = 1\)；
\item 是否存在常数 \(\lambda\)，使得 \(|AB| + |CD| = \lambda |AB| \cdot |CD|\) 恒成立？若存在，求 \(\lambda\) 的值；若不存在，请说明理由.
\end{enumerate}

\bitmapfigure[max width=0.42\linewidth]{img/2010/shandong_science/q21_fig1.png}
\end{problem}

\begin{solution}
\begin{enumerate}
\item 设椭圆的半焦距为 \(c\). 由离心率得
\(\frac ca=\frac{\sqrt2}{2}\)，\(c=\frac{a}{\sqrt2}\)，\(b^2=a^2-c^2=\frac{a^2}{2}\).
椭圆上任一点到两焦点的距离之和为 \(2a\)，而两焦点间的距离为 \(2c\)，所以题设三角形的周长为 \(2a+2c\). 因而
\[
2a+2c=4(\sqrt2+1)
\]
代入 \(c=\frac{a}{\sqrt2}\) 得 \(a=2\sqrt2\)，从而 \(b=2\)，\(c=2\). 椭圆的标准方程为
\[
\frac{x^2}{8}+\frac{y^2}{4}=1
\]
焦点为 \(F_1(-2,0)\)，\(F_2(2,0)\). 等轴双曲线的顶点为这两个焦点，所以其标准方程为
\[
\frac{x^2}{4}-\frac{y^2}{4}=1
\]
\item 设 \(P(u,v)\). 因为 \(P\) 在双曲线上且不是顶点，所以 \(u^2-v^2=4\)，且 \(u\ne\pm2\). 因而
\(k_1=\frac{v}{u+2}\)，\(k_2=\frac{v}{u-2}\)，\(k_1k_2=\frac{v^2}{u^2-4}=1\).
\item 设过焦点 \(F_1(-2,0)\) 的直线斜率为 \(k\)，其方程为 \(y=k(x+2)\). 代入椭圆方程，得
\[
(1+2k^2)x^2+8k^2x+8k^2-8=0
\]
关于 \(x\) 的一元二次方程的判别式为
\[
\Delta=(8k^2)^2-4(1+2k^2)(8k^2-8)=32(1+k^2)
\]
两交点的横坐标之差的绝对值为 \(\frac{\sqrt\Delta}{1+2k^2}\)，而直线斜率为 \(k\)，所以弦长为
\[
L(k)=\frac{\sqrt\Delta}{1+2k^2}\sqrt{1+k^2}=\frac{4\sqrt2(1+k^2)}{1+2k^2}
\]
过另一个焦点 \(F_2(2,0)\) 的直线方程为 \(y=k(x-2)\). 代入椭圆方程得
\[
(1+2k^2)x^2-8k^2x+8k^2-8=0
\]
其判别式仍为 \(32(1+k^2)\)，故弦长仍为
\[
L(k)=\frac{4\sqrt2(1+k^2)}{1+2k^2}
\]
由于 \(P\) 不是双曲线顶点，故 \(v\ne0\)，从而 \(k_1\ne0\). 令 \(t=k_1^2>0\)，由 \(k_1k_2=1\) 得 \(k_2^2=\frac1t\)，于是
\(|AB|=\frac{4\sqrt2(t+1)}{2t+1}\)，\(|CD|=\frac{4\sqrt2(t+1)}{t+2}\).
从而
\(|AB|+|CD|=\frac{12\sqrt2(t+1)^2}{(2t+1)(t+2)}\)，\(|AB|\cdot|CD|=\frac{32(t+1)^2}{(2t+1)(t+2)}\).
二者之比为常数 \(\frac{3\sqrt2}{8}\)，所以存在这样的常数，且
\[
\lambda=\frac{3\sqrt2}{8}
\]
\end{enumerate}
\end{solution}

\begin{problem}
已知函数 \( f(x) = \ln x - ax + \frac{1 - a}{x} - 1 \) （\(a \in \R\)）.
\begin{enumerate}
\item 当 \(a \leqslant \frac{1}{2}\) 时，讨论 \(f(x)\) 的单调性；
\item 设 \( g(x) = x^{2} - 2bx + 4 \)，当 \(a = \frac{1}{4}\) 时，若对任意 \(x_1 \in (0, 2)\)，存在 \(x_2 \in [1, 2]\)，使 \(f(x_1) \geqslant g(x_2)\)，求实数 \(b\) 的取值范围.
\end{enumerate}
\end{problem}

\begin{solution}
\begin{enumerate}
\item 函数定义域为 \(x>0\)，且
\[
f'(x)=\frac1x-a-\frac{1-a}{x^2}=-\frac{(x-1)(ax+a-1)}{x^2}
\]

当 \(a\leqslant0\) 时，\(ax+a-1=a(x+1)-1<0\). 因此当 \(0<x<1\) 时 \(f'(x)<0\)，当 \(x>1\) 时 \(f'(x)>0\)，故 \(f(x)\) 在 \((0,1)\) 上单调递减，在 \((1,+\infty)\) 上单调递增.

当 \(0<a<\frac12\) 时，令 \(r=\frac{1-a}{a}\)，则 \(r>1\)，且
\[
f'(x)=-\frac{a(x-1)(x-r)}{x^2}
\]
所以 \(f(x)\) 在 \((0,1)\) 上单调递减，在 \((1,r)\) 上单调递增，在 \((r,+\infty)\) 上单调递减.

当 \(a=\frac12\) 时，
\[
f'(x)=-\frac{(x-1)^2}{2x^2}\leqslant0
\]
故 \(f(x)\) 在 \((0,+\infty)\) 上单调递减.
\item 当 \(a=\frac14\) 时，
\[
f'(x)=-\frac{(x-1)(x-3)}{4x^2}
\]
所以 \(f(x)\) 在 \((0,1)\) 上递减，在 \((1,2)\) 上递增，故对任意 \(x_1\in(0,2)\)，有
\[
f(x_1)\geqslant f(1)=\ln1-\frac14+\frac34-1=-\frac12
\]

题设条件“对任意 \(x_1\in(0,2)\)，存在 \(x_2\in[1,2]\)，使 \(f(x_1)\geqslant g(x_2)\)”等价于
\[
\min_{1\leqslant x_2\leqslant2}g(x_2)\leqslant-\frac12
\]
事实上，必要性由取 \(x_1=1\) 得到；反之若上式成立，取使 \(g\) 取得最小值的 \(x_2\)，则对任意 \(x_1\) 都有 \(g(x_2)\leqslant-\frac12\leqslant f(x_1)\).

又
\[
g(x_2)\leqslant-\frac12\iff x_2^2-2bx_2+4\leqslant-\frac12\iff b\geqslant \frac{x_2}{2}+\frac{9}{4x_2}
\]
令 \(h(x)=\frac{x}{2}+\frac{9}{4x}\)，则在 \([1,2]\) 上
\[
h'(x)=\frac12-\frac{9}{4x^2}<0
\]
所以 \(h(x)\) 的最小值为 \(h(2)=1+\frac98=\frac{17}{8}\). 因而 \(b\) 的取值范围为
\[
b\geqslant\frac{17}{8}
\]
\end{enumerate}
\end{solution}
